\chapter{暑期试讲与平台改进}

初赛材料完成后，VeriSpecOSLab 在暑假期间面向华东师范大学计算机拔尖班的部分学生进行了两节课试讲。我们将整套工作流带入真实课堂，观察学生能否顺利使用平台、理解操作系统问题并完成早期开发任务。试讲集中暴露了三类问题。平台的安装与操作还不够方便，学生对课程所需前置知识的了解不够充分，独立进行大型项目开发的经验也不足。我们随后围绕这三类问题完善课程材料、学生流程和平台实现。

本章说明试讲前的平台形态、三类课堂问题和对应修改。当前版本的课程入口、渐进式教学、Agent 协作与工程验收均由这次真实教学反馈推动形成。

\section{试讲发现的三类问题}

试讲前的平台已经具备规格、Agent、运行验证和过程记录等基本能力。课程以从零构建个性化操作系统为目标。学生需要在课程开始阶段确定系统目标、参考方案、内核组织和目标平台，随后为不同模块记录职责、接口、操作行为和验收要求。平台按照实验进度开放问答、设计、实现和调试等 AI 功能，并同时提供命令行与网页入口。这套方案覆盖了从设计、实现到验证的完整流程，却在两节试讲中表现出较高的起步门槛。

第一类问题是\textbf{使用不便}。学生进入课程后要处理依赖安装、命令行入口、模型服务、知识库、项目初始化和本地工具链等多项配置。早期命令与概念较多，同一任务又分布在不同入口。本地环境出现问题后，学生往往需要花费较多时间调试，挤占了理解课程内容和完成实验的时间。

第二类问题是\textbf{对前置知识了解不足}。即使学生已经学习 C 语言和计算机组成原理，此前也未必系统了解过操作系统，大部分人对操作系统的认识仍停留在日常使用层面。早期课程要求学生在接触具体机制前决定内核组织、保护模型和资源接口，还要用结构化文档说明系统目标、方案选择和模块行为。这导致学生没有办法理解这些选择将如何影响后续实现。

第三类问题是\textbf{大型项目开发经验不足}。从零构建操作系统需要同时使用 Git、交叉编译器、链接脚本、Make、QEMU、串口日志和调试工具。学生还要阅读 YAML 规格、判断构建错误并设计测试。模型能够快速生成代码，但缺少大型项目开发经验的学生很难对工具给出的结果作出合理评价。试讲表明，平台需要直接支持这些工程步骤，并给学生保留理解和复查结果的入口。

\section{降低平台使用成本}

试讲时，学生要先在 DevBox 与本地安装之间选择。前者需要准备容器环境，后者要分别安装并检查交叉编译器、QEMU、GDB、Bun 和 VOS。每次运行命令还要进入 VOS workspace，并重复指定学生项目路径。若使用网页入口，学生还要分别启动 Agent 后端和 Portal。正式进入实验前，学生已经花费较多时间选择环境、启动服务和检查工具。

早期命令还直接暴露了阶段查询、工具链检查、规格一致性检查、架构组合、Agent 上下文准备、实现计划和代码生成等内部步骤。学生不仅要理解当前 Lab，还要记住这些命令的顺序。遗漏一步，后续流程便无法继续。

试讲中，部分学生在下载或更新 VOS 时受到网络状况影响，安装过程容易中断。改造后，学生入口统一为本地 CLI。仓库依赖安装完成后，学生只需在 CLI 目录执行一次 \code{bun link}，即可把当前仓库中的代码链接为 \code{vos} 命令，无需再从远程地址下载 CLI 或等待运行时更新。进入项目目录后，学生可以直接使用 \code{vos}，也无需启动 Portal、Web UI 和 worker、切换命令行与网页入口或重复输入 workspace 命令和项目路径。

原来分散的准备与执行步骤也被合并。\code{vos init} 创建项目基础文件，\code{vos doctor} 集中检查编译器、构建工具和运行环境，\code{vos agent config} 配置模型服务。\code{spec lint} 统一检查规格，\code{agent implement} 内部完成 Agent 的上下文准备、实现计划和代码生成。学生不再操作阶段管理、架构组合和 Agent 内部流程。

课程材料同时补充前置条件、预计时间、预期结果和自检问题。学生需要记忆的命令更少，平台服务的启动步骤被取消，本地环境问题也有了统一检查入口。

\section{补充课程所需前置知识}

试讲后，课程先帮助学生建立操作系统直觉，再随 Lab 深入具体机制。Lab 1 安排了读取隐藏字符串（flag）的热身任务。教师提供包含两个 flag 文件的文件系统镜像。学生先在 Linux 中编写 C 程序，通过文件接口读取并输出 flag，再在裸机环境中读取同一镜像，自行读取磁盘块、查找目录和文件，并从串口输出相同结果。通过对比两份程序，学生可以看到 Linux 内核提供了哪些文件系统与设备访问能力，又有哪些工作需要裸机程序自行完成。

语言、ISA 和项目身份在课程开始阶段确定。内核组织、保护模型、进程、文件系统和资源接口等选择移到对应 Lab。学生在学习启动、内存、中断和进程等机制后，再逐步更新 DesignSpec。这样，规格记录的是学生已经接触并理解的设计问题，不要求学生预先猜测完整内核的结构。

试讲后，我们重新编写了 Book 与 Lab 两类实验手册。Book 按课程进度解释当前 Lab 所需的概念和代码背景，完整架构的讨论分散到对应阶段。Lab 增加概览、前置依赖、预计时间、任务步骤、预期结果和自检点，并用程序运行结果和 QEMU 输出解释抽象机制。学生先从 Book 理解问题，再跟随 Lab 完成实现和验收，减少在概念说明、零散命令和平台文档之间反复查找的负担。问答助手补充概念解释和方案比较，知识库回答保留资料引用。这些调整降低了实验手册的阅读与理解门槛。

\section{为大型项目开发提供工程支持}

完成这些改造后，学生可以把从零构建内核的长期任务拆到当前 Lab 和模块。DesignSpec 记录系统设计，ModuleSpec 记录模块职责、实现范围、性质和检查项，InterfaceSpec 记录跨模块接口，GoalSpec 记录可度量的扩展目标，SpecPatch 记录跨模块修改的原因和影响。学生每次只需处理与当前阶段直接相关的内容，同时能够从规格中查清模块边界和跨模块影响。

规格仍由学生亲自编写。\code{spec lint} 和只读评审助手负责发现缺少的字段、引用与验收项。规格完善并提交后，实现助手再生成当前模块的代码和测试。整个过程以学生的设计判断为起点，再由助手完成实现，学生不必独自面对空白的内核工程。

进入实现阶段后，助手可以反复尝试，同时不影响学生正在使用的代码目录。助手在单独的临时 Git 工作目录中修改代码，平台从 Git 读取实际变化，并根据 ModuleSpec 和 SpecPatch 检查修改范围。字段错误、修改越界或测试失败会返回同一模型会话继续修正。检查全部通过后，学生可以查看代码差异、测试结果和提交记录。

进入验证阶段后，平台实际运行构建、公开测试、接口契约测试、固定随机种子的模糊测试和有限执行轨迹测试。QEMU Runner 根据串口输出区分正常启动、内核报错和超时。报告把运行结果与 Git 提交、规格摘要和运行配置关联起来。学生可以从报告回查代码、命令和日志，据此评价 AI 给出的实现是否可信，并逐步积累构建、运行、调试和验证内核的工程经验。
